\begin{figure}[!htbp]
\centering
\begin{adjustbox}{max width=\linewidth}
\begin{tikzpicture}[ent/.style={ink, rectangle, fill=fslate, minimum width=22mm, minimum height=8mm, font=\small\bfseries},
    att/.style={ink, ellipse, fill=fsand, minimum height=6.5mm, font=\scriptsize, inner sep=1.5pt},
    rel/.style={ink, diamond, fill=fsage, aspect=1.9, inner sep=1pt, font=\footnotesize\bfseries}]
  \node[ent] (S) at (0,0) {STUDENT};
  \node[rel] (E) at (3.6,0) {ENROLLS};
  \node[ent] (C) at (7.2,0) {COURSE};
  \node[rel] (O) at (10.6,0) {OFFERS};
  \node[ent] (D) at (13.6,0) {DEPARTMENT};
  \draw[thinline] (S) -- node[above, font=\scriptsize] {M} (E);
  \draw[thinline] (E) -- node[above, font=\scriptsize] {N} (C);
  \draw[thinline, double, double distance=1.3pt] (C) -- node[above, font=\scriptsize] {N} (O);
  \draw[thinline] (O) -- node[above, font=\scriptsize] {1} (D);
  \node[att] (s1) at (-1.3,1.4) {\underline{roll\_no}}; \node[att] (s2) at (0.2,1.6) {name}; \node[att] (s3) at (1.5,1.3) {dob};
  \foreach \a in {s1,s2,s3} \draw[thinline] (S) -- (\a);
  \node[att] (e1) at (3.6,1.4) {grade}; \draw[thinline] (E) -- (e1);
  \node[att] (c1) at (6.0,1.4) {\underline{course\_id}}; \node[att] (c2) at (7.4,1.6) {title}; \node[att] (c3) at (8.6,1.3) {credits};
  \foreach \a in {c1,c2,c3} \draw[thinline] (C) -- (\a);
  \node[att] (d1) at (13.6,1.4) {\underline{dept\_id}}; \draw[thinline] (D) -- (d1);
\end{tikzpicture}
\end{adjustbox}
\caption{An ER diagram for a university. ENROLLS is many-to-many and carries the attribute \emph{grade};
every course must be offered by a department (total participation, double line).}
\end{figure}
